\documentclass[12pt]{article}
\usepackage[utf8]{inputenc}
\usepackage[T1]{fontenc}
\usepackage{hyperref}
\title{Yandruwandha Language Resources for K-12 Curriculum Development}
\author{Compiled by Claude Sonnet 5.5 (claude-sonnet-5-5) for Andr\'as Kornai}
\date{}
\begin{document}
\maketitle

\section{Introduction}
Yandruwandha (ISO 639-3 ynd; Glottocode yand1253) is an extinct Pama-Nyungan language of Australia, associated with the Innamincka area. No living first-language speakers are known, so any school curriculum would be a heritage and revitalization programme built on archival documentation. The search found one main body of material, the work of Gavan Breen, published in 2015. I found no school programme, primer, or audio collection in the search results.

\section{Primary Educational Resources}
No primers, textbooks or school programmes were found. The dictionary and stories below are the closest material to a classroom resource. They were written with descendants of the original people of the area in mind, and the stories could be adapted for teaching.
\begin{itemize}
\item Innamincka Words: Yandruwandha dictionary and stories (ANU Press, 2015): \url{https://press.anu.edu.au/publications/innamincka-words-yandruwandha-dictionary-and-stories}
\end{itemize}

\section{Linguistic and Reference Materials}
\begin{itemize}
\item Innamincka Words: open-access copy in the OAPEN library: \url{https://library.oapen.org/handle/20.500.12657/32427}
\item Full-text PDF in OAPEN (the search result did not say which volume this PDF is): \url{https://library.oapen.org/bitstream/id/89c1553e-930f-4f3b-ab11-034bf6b1f339/610991.pdf}
\item Wikipedia overview of the language: \url{https://en.wikipedia.org/wiki/Yandruwandha_language}
\end{itemize}
A companion grammar, \emph{Innamincka Talk: a grammar of the Innamincka dialect of Yandruwandha with notes on other dialects} (Breen, ANU Press, 2015), was named in the search results. No direct URL for it appeared in them, so none is given here. It is technical and aimed at specialists.

\section{Religious and Community Resources}
No scripture translations or community organisations were found in the search results.

\section{Digital and Audio Resources}
No audio, video or software was found. The digital resources are the open-access ANU Press and OAPEN editions above.

\section{Curriculum Development Context}
Because the language is extinct, a curriculum would have to be reconstructed from the Breen dictionary, stories and grammar. This would ideally be done with descendant communities, whose cultural authority over the material should be respected. Nothing found indicates that any school currently teaches the language.
\end{document}